%!TEX root = main.tex

%\section{FIRRTL Syntax}

%%!TEX root = main.tex

\begin{figure}[htbp]
{\footnotesize
\[
\begin{array}{l l l}
%\mbox{\bf Circuit} & 
circuit & = & \mbox{``circuit''}, id, newline, indent, \{module \mid extmodule\}, dedent;\\
%
%\mbox{\bf Module} & 
module & = & \mbox{``module''} \ id, newline, indent, \\
%& 
& & \{port, newline\}, \{stmt, newline\}, dedent;\\
%
%\mbox{\bf External Module} & 
extmodule & = & \mbox{``extmodule''}, id, newline, indent,\\
%& 
  &  &  \{port, newline\}, [\mbox{``defname''}, \mbox{``=''}, id, newline], \\
% & 
 & & \{\mbox{``parameter''}, \mbox{``=''}, (string \mid int), newline\}, dedent;\\
%
%\mbox{\bf Port} & 
port & = & (\mbox{``input''} \mid \mbox{``output''}), id, \mbox{``:''}, type;\\
%
 %\mbox{\bf Statment} & 
 stmt & = & \mbox{``wire''}, id, \mbox{``:''}, type \\
 %& 
 & &  \mid \mbox{``reg''}, id, \mbox{``:''}, type, \mbox{``,''}, exp, [\mbox{``(with:\{reset=>(''}, exp, \mbox{``,''}, exp, \mbox{``)\})''}] \\
 %& 
  &   &  \mid memory \mid \mbox{``inst''}, id, \mbox{``of''}, id \mid \mbox{``node''}, id, \mbox{``=''}, id  \\
 %& 
  &   & \mid reference, \mbox{``<=''}, exp  \mid reference, \mbox{``<-''}, exp  \\
 %& 
  &   &\mid reference, \mbox{``is invalid''}  \mid \mbox{``attach(''}, \{reference\}, \mbox{``)''}  \\ 
%&  
&   &\mid \mbox{``when''}, exp, \mbox{``:''}, newline, indent, \{stmt\}, dedent \\
%& 
& & [\mbox{``else:''}, indent, \{stmt\}, dedent] \\ 
%&  
&   &\mid \mbox{``stop(''}, exp, \mbox{``,''}, exp, \mbox{``,''}, int, \mbox{``)''}\\ 
%&  
&   &\mid \mbox{``printf(''}, exp, \mbox{``,''}, exp, \mbox{``,''}, string, \mbox{``,''}, \{exp\}, \mbox{``)''}, [\mbox{``:''}, id]\\ 
%&  
&   &\mid \mbox{``skip''};\\ 
%
%\mbox{\bf Memory} & 
memory & = &  \mbox{``mem''}, id, newline, indent, \\
 %& 
 & & \mbox{``data-type''}, \mbox{``=>''}, type, newline, \\
 %& 
 & & \mbox{``depth''}, \mbox{``=>''}, int, newline, \\
 %& 
 & & \mbox{``read-latency''}, \mbox{``=>''}, int, newline, \\ 
 %& 
 & & \mbox{``write-latency}, \mbox{``=>''}, int, newline,  \\
 %& 
 & & \mbox{``read-under-write''}, \mbox{``=>''}, ruw, newline, \\
  %& 
  & & \{\mbox{``reader''}, \mbox{``=>''}, id, newline\}, \\
  %& 
  & & \{\mbox{``writer''}, \mbox{``=>''}, id, newline\}, \\
  %& 
  & & \{\mbox{``readwriter''}, \mbox{``=>''}, id, newline\}, dedent;\\    
%
%\mbox{\bf Read Under Write Flag} & 
ruw & = & \mbox{``old''} \mid \mbox{``new''} \mid \mbox{``undefined''};\\
%
%\mbox{\bf Expression} & 
exp & = & (\mbox{``UInt''} \mid \mbox{``SInt''}), [width], \mbox{``(''}, int, \mbox{``)''} \mid reference \\
%& 
& & \mid  \mbox{``mux''}, \mbox{``(''}, exp, \mbox{``,''} exp, \mbox{``,''}, exp, \mbox{``)''} \mid \mbox{``validif''}, \mbox{``(''}, exp, \mbox{``,''}, exp, \mbox{``)''}\\
%& 
& & \mid primop;\\
%
%& 
reference & = & id \mid reference, \mbox{``.''}, id \mid reference, \mbox{``[''}, int, \mbox{``]''} \mid reference, \mbox{``[''}, exp, \mbox{``]''};\\
%
%& 
width & = & \mbox{``<''}, int, \mbox{``>''};\\
%
%& 
binarypoint & = & \mbox{``<''}, \mbox{``<''}, int, \mbox{``>''}, \mbox{``>''};\\
%
 %\mbox{\bf Bundle Field} & 
 field & = & [\mbox{``flip''}]\ id: type;\\
%
%\mbox{\bf Data Types} & 
type & = & type\_ground \mid type\_aggregate; \\
 %& 
 type\_ground & = &  \mbox{``Clock''} \mid \mbox{``Reset''} \mid \mbox{``AsyncReset''} \mid (\mbox{``UInt''} \mid \mbox{``SInt''}  \mid \mbox{``Analog''}), [width]  \\
 %& 
 & & \mid \mbox{``Fixed''}, [width], [binarypoint];\\
 %& 
  type\_aggregate & = & \mbox{``\{''}, field, \{field\}, \mbox{``\}''} \mid type, \mbox{``[''}, int, \mbox{``]''};\\
%& & & \mbox{``UInt''}, [\mbox{``<''}, int, \mbox{``>''}] \mid \mbox{``SInt''}, [\mbox{``<''}, int, \mbox{``>''} ] \\
%& & & \mid \mbox{``Fixed''}, [\mbox{``<''}, int, \mbox{``>''} ] [\mbox{``<}\mbox{<''}, int, \mbox{``>}\mbox{>''}]\\
% &  &  & \mid \mbox{``Clock''} \mid \mbox{``Analog''},  [\mbox{``<''}, int, \mbox{``>''}] \mid \mbox{``\{''}, field, \mbox{``\}''} \mid type, \mbox{``[''}, int, \mbox{``]''}; \\
%
%\mbox{\bf Primitive Operation} & 
primop & = & primop\_2exp \mid primop\_1exp \mid primop\_1exp1int \mid primop\_1exp2int; \\
%& 
primop\_2exp & = & primop\_2exp\_keyword, \mbox{``(''} , exp, \mbox{``,''} , exp, \mbox{``)''}; \\
%
%& 
primop\_2exp\_keyword & = &  \mbox{``add''}  \mid  \mbox{``sub''}  \mid  \mbox{``sub''} \mid \mbox{``mul''} \mid \mbox{``div''} \mid \mbox{``mod''} \mid \mbox{``lt''} \mid \mbox{``leq''}  \\
%& 
& &  \mid \mbox{``gt''} \mid \mbox{``geq''} \mid \mbox{``eq''} \mid \mbox{``neq''} \mid \mbox{``dshl''}  \mid \mbox{``dshr''} \\
%& 
 & & \mid \mbox{``and''} \mid \mbox{``or''} \mid \mbox{``xor''} \mid \mbox{``cat''}; \\
 %
%& 
primop\_1exp & = & primop\_1exp\_keyword, \mbox{``(''}, exp, \mbox{``)''}; \\
%& 
primop\_1exp\_keyword & = & \mbox{``asUInt''} \mid \mbox{``asSInt''} \mid \mbox{``asClock''} \mid \mbox{``cvt''} \mid \mbox{``neg''} \mid \mbox{``not''}  \\
%& 
& & \mid \mbox{``andr''} \mid \mbox{``orr''} \mid  \mbox{``xorr''};\\
%
%& 
primop\_1exp1int & = & primop\_1exp1int\_keyword, \mbox{``(''}, exp, \mbox{``,''}, int, \mbox{``)''}; \\
%& 
primop\_1exp1int\_keyword & = & \mbox{``pad''} \mid \mbox{``shl''} \mid \mbox{``shr''} \mid \mbox{``head''} \mid \mbox{``tail''}; \\
%& 
primop\_1exp2int & = & primop\_1exp2int\_keyword, \mbox{``(''}, exp, \mbox{``,''}, int, \mbox{``,''}, int \mbox{``)''}; \\
%& 
primop\_1exp2int\_keyword & = & \mbox{``bits''}; \\
%
%\mbox{\bf File Information Token} & 
%info & = & \mbox{``@''}, \mbox{``[''}, \{string, \mbox{`` ''}, linecol\}, \mbox{``]''};\\
%& 
linecol & = &  digit\_dec, \{digit\_dec\}, \mbox{``:''}, digit\_dec, \{digit\_dec\};\\
%& 
id & = & ( \mbox{``\_''} \mid letter), \{ \mbox{``\_''} \mid letter \mid digit\_dec\};
%
\end{array}
\]
}
\caption{Syntax of FIRRTL}
\label{fig-firrtl-synt}
\end{figure}



%%%%%%%%%%%%%%%%%%%%%%%%%%%%%%%%%%%%%%%%%%%%%%%
%%%%%%%%%%%%%%%%%%syntax with info %%%%%%%%%%%%%%%%%%%
%%%%%%%%%%%%%%%%%%%%%%%%%%%%%%%%%%%%%%%%%%%%%%%

\hide{

\begin{figure}[t]
{\footnotesize
\[
\begin{array}{l l l}
%\mbox{\bf Circuit} & 
circuit & = & \mbox{``circuit''}, id, \mbox{``:''}, [info], newline, indent, \{module \mid extmodule\}, dedent;\\
%
%\mbox{\bf Module} & 
module & = & \mbox{``module''} \ id, \mbox{``:''}, [info], newline, indent, \\
%& 
& & \{port, newline\}, \{stmt, newline\}, dedent;\\
%
%\mbox{\bf External Module} & 
extmodule & = & \mbox{``extmodule''}, id, \mbox{``:''}, [info], newline, indent,\\
%& 
  &  &  \{port, newline\}, [\mbox{``defname''}, \mbox{``=''}, id, newline], \\
% & 
 & & \{\mbox{``parameter''}, \mbox{``=''}, (string \mid int), newline\}, dedent;\\
%
%\mbox{\bf Port} & 
port & = & (\mbox{``input''} \mid \mbox{``output''}), id, \mbox{``:''}, type, [info];\\
%
 %\mbox{\bf Statment} & 
 stmt & = & \mbox{``wire''}, id, \mbox{``:''}, type, [info] \\
 %& 
 & &  \mid \mbox{``reg''}, id, \mbox{``:''}, type, \mbox{``,''}, exp, [\mbox{``(with:\{reset=>(''}, exp, \mbox{``,''}, exp, \mbox{``)\})''}], [info] \\
 %& 
  &   &  \mid memory \mid \mbox{``inst''}, id, \mbox{``of''}, id, [info] \mid \mbox{``node''}, id, \mbox{``=''}, id, [info]  \\
 %& 
  &   & \mid reference, \mbox{``<=''}, exp, [info]  \mid reference, \mbox{``<-''}, exp, [info]  \\
 %& 
  &   &\mid reference, \mbox{``is invalid''}, [info]  \mid \mbox{``attach(''}, \{reference\}, \mbox{``)''}, [info]  \\ 
%&  
&   &\mid \mbox{``when''}, exp, \mbox{``:''}, [info], newline, indent, \{stmt\}, dedent \\
%& 
& & [\mbox{``else:''}, indent, \{stmt\}, dedent] \\ 
%&  
&   &\mid \mbox{``stop(''}, exp, \mbox{``,''}, exp, \mbox{``,''}, int, \mbox{``)''}, [info]\\ 
%&  
&   &\mid \mbox{``printf(''}, exp, \mbox{``,''}, exp, \mbox{``,''}, string, \mbox{``,''}, \{exp\}, \mbox{``)''}, [\mbox{``:''}, id], [info]\\ 
%&  
&   &\mid \mbox{``skip''}, [info];\\ 
%
%\mbox{\bf Memory} & 
memory & = &  \mid  \mbox{``mem''}, id, \mbox{``:''}, [info], newline, indent, \\
 %& 
 & & \mbox{``data-type''}, \mbox{``=>''}, type, newline, \\
 %& 
 & & \mbox{``depth''}, \mbox{``=>''}, int, newline, \\
 %& 
 & & \mbox{``read-latency''}, \mbox{``=>''}, int, newline, \\ 
 %& 
 & & \mbox{``write-latency}, \mbox{``=>''}, int, newline,  \\
 %& 
 & & \mbox{``read-under-write''}, \mbox{``=>''}, ruw, newline, \\
  %& 
  & & \{\mbox{``reader''}, \mbox{``=>''}, id, newline\}, \\
  %& 
  & & \{\mbox{``writer''}, \mbox{``=>''}, id, newline\}, \\
  %& 
  & & \{\mbox{``readwriter''}, \mbox{``=>''}, id, newline\}, dedent;\\    
%
%\mbox{\bf Read Under Write Flag} & 
ruw & = & \mbox{``old''} \mid \mbox{``new''} \mid \mbox{``undefined''};\\
%
%\mbox{\bf Expression} & 
exp & = & (\mbox{``UInt''} \mid \mbox{``SInt''}), [width], \mbox{``(''}, int, \mbox{``)''} \mid reference \\
%& 
& & \mid  \mbox{``mux''}, \mbox{``(''}, exp, \mbox{``,''} exp, \mbox{``,''}, exp, \mbox{``)''} \mid \mbox{``validif''}, \mbox{``(''}, exp, \mbox{``,''}, exp, \mbox{``)''}\\
%& 
& & \mid primop;\\
%
%& 
reference & = & id \mid reference, \mbox{``.''}, id \mid reference, \mbox{``[''}, int, \mbox{``]''} \mid reference, \mbox{``[''}, exp, \mbox{``]''};\\
%
%& 
width & = & \mbox{``<''}, int, \mbox{``>''};\\
%
%& 
binarypoint & = & \mbox{``<''}, \mbox{``<''}, int, \mbox{``>''}, \mbox{``>''};\\
%
 %\mbox{\bf Bundle Field} & 
 field & = & [\mbox{``flip''}]\ id: type;\\
%
%\mbox{\bf Data Types} & 
type & = & type\_ground \mid type\_aggregate; \\
 %& 
 type\_ground & = &  \mbox{``Clock''} \mid \mbox{``Reset''} \mid \mbox{``AsyncReset''} \mid (\mbox{``UInt''} \mid \mbox{``SInt''}  \mid \mbox{``Analog''}), [width]  \\
 %& 
 & & \mid \mbox{``Fixed''}, [width], [binarypoint];\\
 %& 
  type\_aggregate & = & \mbox{``\{''}, field, \{field\}, \mbox{``\}''} \mid type, \mbox{``[''}, int, \mbox{``]''};\\
%& & & \mbox{``UInt''}, [\mbox{``<''}, int, \mbox{``>''}] \mid \mbox{``SInt''}, [\mbox{``<''}, int, \mbox{``>''} ] \\
%& & & \mid \mbox{``Fixed''}, [\mbox{``<''}, int, \mbox{``>''} ] [\mbox{``<}\mbox{<''}, int, \mbox{``>}\mbox{>''}]\\
% &  &  & \mid \mbox{``Clock''} \mid \mbox{``Analog''},  [\mbox{``<''}, int, \mbox{``>''}] \mid \mbox{``\{''}, field, \mbox{``\}''} \mid type, \mbox{``[''}, int, \mbox{``]''}; \\
%
%\mbox{\bf Primitive Operation} & 
primop & = & primop\_2exp \mid primop\_1exp \mid primop\_1exp1int \mid primop\_1exp2int; \\
%& 
primop\_2exp & = & primop\_2exp\_keyword, \mbox{``(''} , exp, \mbox{``,''} , exp, \mbox{``)''}; \\
%
%& 
primop\_2exp\_keyword & = &  \mbox{``add''}  \mid  \mbox{``sub''}  \mid  \mbox{``sub''} \mid \mbox{``mul''} \mid \mbox{``div''} \mid \mbox{``mod''} \mid \mbox{``lt''} \mid \mbox{``leq''}  \\
%& 
& &  \mid \mbox{``gt''} \mid \mbox{``geq''} \mid \mbox{``eq''} \mid \mbox{``neq''} \mid \mbox{``dshl''}  \mid \mbox{``dshr''} \\
%& 
 & & \mid \mbox{``and''} \mid \mbox{``or''} \mid \mbox{``xor''} \mid \mbox{``cat''}; \\
 %
%& 
primop\_1exp & = & primop\_1exp\_keyword, \mbox{``(''}, exp, \mbox{``)''}; \\
%& 
primop\_1exp\_keyword & = & \mbox{``asUInt''} \mid \mbox{``asSInt''} \mid \mbox{``asClock''} \mid \mbox{``cvt''} \mid \mbox{``neg''} \mid \mbox{``not''}  \\
%& 
& & \mid \mbox{``andr''} \mid \mbox{``orr''} \mid  \mbox{``xorr''};\\
%
%& 
primop\_1exp1int & = & primop\_1exp1int\_keyword, \mbox{``(''}, exp, \mbox{``,''}, int, \mbox{``)''}; \\
%& 
primop\_1exp1int\_keyword & = & \mbox{``pad''} \mid \mbox{``shl''} \mid \mbox{``shr''} \mid \mbox{``head''} \mid \mbox{``tail''}; \\
%& 
primop\_1exp2int & = & primop\_1exp2int\_keyword, \mbox{``(''}, exp, \mbox{``,''}, int, \mbox{``,''}, int \mbox{``)''}; \\
%& 
primop\_1exp2int\_keyword & = & \mbox{``bits''}; \\
%
%\mbox{\bf File Information Token} & 
info & = & \mbox{``@''}, \mbox{``[''}, \{string, \mbox{`` ''}, linecol\}, \mbox{``]''};\\
%& 
linecol & = &  digit\_dec, \{digit\_dec\}, \mbox{``:''}, digit\_dec, \{digit\_dec\};\\
%& 
id & = & ( \mbox{``\_''} \mid letter), \{ \mbox{``\_''} \mid letter \mid digit\_dec\};
%
\end{array}
\]
}
\caption{Syntax of FIRRTL}
\label{fig-firrtl-synt}
\end{figure}


}

\section{Missing details in Section~\ref{sec-lofirrtl-sem}}\label{app-lofirrtl}

\paragraph*{The definitions of $te(e)$ and $s(e)$}
%
For $\theta \bwidth{n}(n')$ such that $\theta \in \{\mbox{\prettyll{UInt}, \prettyll{SInt}}\}$, $n > 0$, and $n'$ is an integer, we define $bv_{\theta \bwidth{n}(n')}$ as follows.
\begin{itemize}
\item For $n > 0$ and $n' \ge 0$ such that $n' < 2^n$, let $bv_{UInt<n>(n')}$ denote the bit vector of length $n$ representing the unsigned integer $n'$.
%
\item For $n > 0$ and an integer $n'$ such that $-2^{n-1} \le n' < 2^{n-1}$, let $bv_{SInt<n>(n')}$ denote the bit vector of length $n$ representing the signed integer $n'$.
\end{itemize}
The function $te$ and $s$ from $\cpnts(A)$ to $\gtypes$ are extended to expressions as follows.
\begin{itemize}
\item $te(\theta \bwidth{n}(c)) = \theta \bwidth{n}$ and $s(\theta \bwidth{n}(c)) = bv_{\theta \bwidth{n}(c)}$, where $\theta \in \{UInt, SInt\}$,
%\item $te(\sint{n}(c)) = \sint{n}$,
%
\item let $te(e_1) = \theta \bwidth{n_1}$ and $te(e_2) = \theta \bwidth{n_2}$, where $\theta \in \{UInt, SInt\}$, then $te(mux(e_0, e_1, e_2)) = \theta\bwidth{\max(n_1, n_2)}$, moreover, if $s(e_0) = 1$, then $s(mux(e_0, e_1, e_2)) = bv_{\theta\bwidth{n_0}(s(e_1))}$, otherwise, $s(mux(e_0, e_1, e_2)) = bv_{\theta\bwidth{n_0}(s(e_2))}$,
%
\item $te(validif(e_1, e_2)) = te(e_2)$, moreover, if $s(e_1) = 1$, then $s(validif(e_1, e_2)) = s(e_2)$, otherwise, $s(validif(e_1, e_2))  = \bot$,
%
\item let $te(e_1) = \theta \bwidth{n_1}$ and $te(e_2) = \theta \bwidth{n_2}$, where $\theta \in \{UInt, SInt\}$, then
\begin{itemize}
\item $te(add(e_1, e_2)) =\theta \bwidth{\max(n_1, n_2)+1}$ and $te(sub(e_1, e_2)) = \theta \bwidth{\max(n_1, n_2)+1}$, moreover, $s(add(e_1, e_2)) = bv_{\theta \bwidth{\max(n_1,n_2)+1} (s(e_1) + s(e_2))}$ and $s(sub(e_1, e_2)) = bv_{\theta \bwidth{\max(n_1,n_2)+1} (s(e_1) - s(e_2))}$,
%
\item $te(mul(e_1, e_2)) = \theta \bwidth{n_1+n_2}$ and $s(mul(e_1, e_2)) = bv_{\theta \bwidth{n_1+n_2} (s(e_1) \times s(e_2))}$,
%
\item  if $\theta = UInt$, then $te(div(e_1, e_2)) = UInt \bwidth{n_1}$ and $s(div(e_1, e_2)) = bv_{UInt \bwidth{n_1} (s(e_1) \div s(e_2))}$, otherwise, $te(div(e_1, e_2)) = SInt \bwidth{n_1+1}$ and $s(div(e_1, e_2)) = bv_{SInt \bwidth{n_1+1} (s(e_1) \div s(e_2))}$,
%
\item $te(mod(e_1, e_2)) = \theta \bwidth{\min(n_1,n_2)}$ and $s(mod(e_1, e_2)) = bv_{\theta \bwidth{\min(n_1,n_2)} (s(e_1) \bmod s(e_2))}$,
%
\item $te(cp(e_1, e_2)) = UInt \bwidth{1}$, moreover, $s(cp(e_1, e_2)) = 1$ iff $s(e_1)\ cp\ s(e_2)$, where $cp \in \{lt, leq, gt, geq, eq, neq\}$,
%
\item $te(dshl(e_1, e_2)) = \theta \bwidth{n_1+2^{n_2}-1}$ and $s(dshl(e_1, e_2)) = bv_{\theta \bwidth{n_1+2^{n_2}-1}(s(e_1) \times 2^{s(e_2)})}$,
%
\item $te(dshr(e_1, e_2)) = \theta \bwidth{n_1}$ and $s(dshr(e_1, e_2)) = bv_{\theta \bwidth{n_1}(s(e_1) \div 2^{s(e_2)})}$,
%
\item $te(and(e_1, e_2)) = te(or(e_1, e_2)) =  te(xor(e_1, e_2)) = UInt \bwidth{\max(n_1, n_2)}$, moreover,
$s(and(e_1, e_2))$ is the bitwise and of $bv_{\theta\bwidth{\max(n_1, n_2)}(s(e_1))}$ and $bv_{\theta\bwidth{\max(n_1, n_2)}(s(e_2))}$,  $s(or(e_1, e_2))$ is the bitwise or of $bv_{\theta\bwidth{\max(n_1, n_2)}(s(e_1))}$ and $bv_{\theta\bwidth{\max(n_1, n_2)}(s(e_2))}$, $s(xor(e_1, e_2))$ is the bitwise xor of $bv_{\theta\bwidth{\max(n_1, n_2)}(s(e_1))}$ and $bv_{\theta\bwidth{\max(n_1, n_2)}(s(e_2))}$,
%
\item $te(cat(e_1, e_2)) = UInt \bwidth{n_1 + n_2}$, moreover,
\begin{itemize}
\item if either $\theta = UInt$, or $\theta=SInt$, $s(e_1) \ge 0$, and $s(e_2) \ge 0$, then $s(cat(e_1, e_2))$ is $bv_{UInt\bwidth{n_1+n_2}(s(e_1) \times 2^{n_2} + s(e_2))}$,
%
\item if  $\theta=SInt$, $s(e_1) \ge 0$, and $s(e_2) < 0$, then $s(cat(e_1, e_2))$ is $bv_{UInt\bwidth{n_1+n_2}(s(e_1) \times 2^{n_2} + 2^{n_2} + s(e_2))}$,
%
\item if  $\theta=SInt$ and $s(e_1) < 0, s(e_2) \ge 0$ $s(cat(e_1, e_2))$ is $bv_{UInt\bwidth{n_1+n_2}((2^{n_1}+s(e_1)) \times 2^{n_2} + s(e_2))}$,
%
\item if  $\theta=SInt$ and $s(e_1) < 0, s(e_2) < 0$ $s(cat(e_1, e_2))$ is $bv_{UInt\bwidth{n_1+n_2}((2^{n_1}-s(e_1)) \times 2^{n_2} + 2^{n_2}-s(e_2))}$,
\end{itemize}
\end{itemize}
%
\item let $te(e_1) = \theta \bwidth{n_1}$ and $te(e_2) = UInt \bwidth{n_2}$, where $\theta \in \{UInt, SInt\}$, then
\begin{itemize}
\item $te(dshl(e_1, e_2)) = \theta \bwidth{n_1+2^{n_2}-1}$ and $s(dshl(e_1, e_2)) = bv_{\theta \bwidth{n_1+2^{n_2}-1}(s(e_1) \times 2^{s(e_2)})}$,
%
\item $te(dshr(e_1, e_2)) = \theta \bwidth{n_1}$ and $s(dshr(e_1, e_2)) = bv_{\theta \bwidth{n_1}(s(e_1) \div 2^{s(e_2)})}$,
%
\end{itemize}
%
\item let $te(e_1) = \theta \bwidth{n_1}$, where $\theta \in \{UInt, SInt\}$, then
\begin{itemize}
\item $te(andr(e_1)) = te(orr(e_1)) =  te(xorr(e_1)) = \uint{1}$, moreover, $s(andr(e_1))$ (resp. $s(orr(e_1))$, $s(xorr(e_1))$) is the and (resp. or, xor) of all the bits in the bit vector $s(e_1)$,
%
\item $te(neg(e_1)) = \sint{n_1 +1}$, $te(not(e_1)) =  \uint{n_1}$, moreover,
\begin{itemize}
\item $s(neg(e_1)) = bv_{SInt\bwidth{n_1+1}(-s(e_1))}$,
%
\item $s(not(e_1))$ is obtained by applying the logical not to each bit of the bit vector $s(e_1)$,
\end{itemize}
%\begin{itemize}
%\item if $\theta = UInt$, then $s(neg(e_1)) = bv_{SInt\bwidth{n_1+1}(-s(e_1))}$,
%
%\item if $\theta = SInt$, then $s(neg(e_1)) = bv_{SInt\bwidth{n_1+1}(-s(e_1))}$.
%\end{itemize}
%
\item $te(cvt(e_1)) = \sint{n_1 +1}$ if $\theta = UInt$, and $te(cvt(e_1)) =  \sint{n_1}$ otherwise,  moreover,
\begin{itemize}
\item if $\theta = UInt$, then $s(cvt(e_1)) = bv_{SInt\bwidth{n_1+1}(s(e_1))}$,
%
\item if $\theta = SInt$, then $s(cvt(e_1)) = s(e_1)$,
\end{itemize}
\end{itemize}
%
\item let $te(e_1) = \theta \bwidth{n_1}$, where $\theta \in \{UInt, SInt\}$, then
\begin{itemize}
\item $te(pad(e_1, n)) = \theta\bwidth{\max(n_1, n)}$, moreover, if $n_1 < n$, then $s(pad(e_1,n)) = bv_{\theta\bwidth{n}(s(e_1))}$, otherwise, $s(pad(e_1,n)) = s(e_1)$,
%
\item $te(shl(e_1,n)) = \theta \bwidth{n_1 + n}$, moreover, $s(shl(e_1,n)) =  bv_{\theta\bwidth{n_1 + n}(2^{n} \times s(e_1))}$,
%\begin{itemize}
%\item if $\theta = UInt$, then $s(shl(e_1,n)) =  bv_{UInt\bwidth{n_1 + n}(2^{n} \times s(e_1))}$,
%\item if $\theta = SInt$, then $s(shl(e_1,n)) = bv_{SInt\bwidth{n_1 + n}(2^{n} \times s(e_1))}$,
%\end{itemize}
%
\item $te(shr(e_1, n)) = \theta \bwidth{n_1-n}$ if $n_1 > n$, and $te(shr(e_1, n)) =  \theta \bwidth{1}$ otherwise,  moreover,
\begin{itemize}
\item if $n_1 > n$, then $s(shr(e_1, n))$ is obtained from $s(e_1)$ by shifting $n$ bits right,
%
\item otherwise, if $\theta = UInt$, then $s(shr(e_1, n))=bv_{UInt<1>(0)}$, otherwise, $s(shr(e_1, n))=bv_{SInt<1>(0)}$ if $s(e_1) \ge 0$, and $s(shr(e_1, n))=bv_{SInt<1>(1)}$ if $s(e_1) < 0$,
\end{itemize}
%
\item $te(head(e_1, n)) = \uint{n}$, $te(tail(e_1, n)) = \uint{n_1-n}$, moreover, $s(head(e_1, n))$ is the most significant $n$ bits of $s(e_1)$, $s(tail(e_1, n))$ is obtained from $s(e_1)$ by truncating the $n$ most significant bits of $s(e_1)$,
%
\item $te(bits(e_1, n, n')) = \uint{n'-n+1}$, and $s(bits(e_1, n, n'))$ is the bit vector comprising the $n, n+1, \cdots, n'$-th bit of $s(e_1)$.
\end{itemize}
%
\end{itemize}


\section{Axiomatic semantics of compilation passes}\label{sec-appen-sem-cpass}

The semantics of compilation pass InferType on a module $M$ is specified by the following formula,
\[
\begin{array}{l c l}
\varphi_{InferType, M}(ce, cm, ce', cm') & \egdef & \exists (ce_i)_{0 \le i \le n} (cm_i)_{0 \le i \le n}.\ ce = ce_0 \wedge cm = cm_0 \wedge ce' = ce_n\ \wedge \\
& &  cm' = cm_n \wedge \bigwedge \limits_{1 \le i \le n} \varphi_{InferType, s_i}(ce_{i-1}, cm_{i-1}, ce_i, cm_i),
\end{array}
\]
where the semantics of $\varphi_{InferType, s_i}$ is defined by the rules in Table~\ref{tab:infertype-sem}. Moreover,
$type_{ce}(e)$ in Table~\ref{tab:infertype-sem} is defined in the same flavor as $te(e)$ in the definition of operational semantics of LoFIRRTL. For instance, $type_{ce}(add(e_1, e_2))$ is defined as follows:
\begin{itemize}
\item if $type_{ce}(e_1) = (\mbox{``Gtype''}, \theta\mbox{<}n_1\mbox{>})$ and  $type_{ce}(e_2) = (\mbox{``Gtype''}, \theta\mbox{<}n_2\mbox{>})$ for some $\theta \in \{UInt, SInt\}$, then $type_{ce}(add(e_1, e_2)) = (\mbox{``Gtype''}, \theta\mbox{<}\max(n_1,n_2)\mbox{>})$,
%\item if $type_{ce}(e_1) = (\mbox{``Gtype''}, \sint{n_1})$ and  $type_{ce}(e_2) = (\mbox{``Gtype''}, \sint{n_2})$, then $type_{ce}(add(e_1, e_2)) = (\mbox{``Gtype''}, \sint{\max(n_1,n_2)})$,
\item otherwise, $type_{ce}(add(e_1, e_2)) = \bot$.
\end{itemize}


\begin{table}[htbp]
  \footnotesize
  \begin{gather*}
%    \infer[skip]
%    {\seqq{\varphi_{InferType, \mbox{skip}}((ce, cm), (ce, cm))}}
%    {\seqq{-}}
%    \quad
    \infer[wire]
%    {\seqq{}}
    {\seqq{\varphi_{InferType, \mbox{wire }v:t}((ce, cm), (ce', cm))}}
   {\seqq{ce' = ce [(\mbox{``AggrType''}, t, \mbox{``Wire''})/v]}}
    \\[1ex]
    \infer[reg]
    {\varphi_{InferType, \mbox{reg } r:t, c \mbox{ with: reset => }(e_1, e_2)}((ce, cm), (ce', cm))}
   {\seqq{ce' = ce [(\mbox{``RegType''}, t, \mbox{``Register''})/r]}}
    \quad
    \\[1ex]
    \infer[node]
    {\seqq{\varphi_{InferType, \mbox{node } v = e} ((ce, cm), (ce', cm'))}}
   {\seqq{ce' = ce[(\mbox{``AggrType''}, type_{ce}(e), \mbox{``Node''})/v]}}
    \quad
    \\[1ex]
    \infer[inst of]
    {\varphi_{InferType, \mbox{inst } v_1 \mbox{ of } v_2} ((ce, cm), (ce', cm'))}
   {\seqq{ce' = ce[(prj_{1,2}(type_{ce}(v_2)), \mbox{``InstanceOf''})/v_1] }}
  \quad
    \\[1ex]
    \infer[$st$: skip, when, connect, invalid statement]
    {\varphi_{InferType, st} ((ce, cm), (ce, cm))}
%   {\seqq{rs(v) \neq \bot}} {\seqq{rs' = rs[s(e)/v]}}
   {\seqq{-}}
  \end{gather*}
%  \vspace{-4mm}
  \caption{Semantics of InferType compilation pass}
  \label{tab:infertype-sem}
\end{table}

%%%%%%%%%%%%%%%%%%%%%%%%%%%%%%%
%%%%%%%%%definition of type_ce(e) omitted %%%%%%
%%%%%%%%%%%%%%%%%%%%%%%%%%%%%%%
\hide{
\begin{table}[htbp]
  \footnotesize
  \begin{gather*}
%    \infer[skip]
%    {\seqq{\varphi_{InferType, \mbox{skip}}((ce, cm), (ce, cm))}}
%    {\seqq{-}}
%    \quad
    \infer[UInt const]
%    {\seqq{}}
    {\seqq{type_{ce}(\uint{n}(c)) = (\mbox{``Gtyp''}, \uint{n})}}
   {\seqq{-}}
   \quad
    \infer[SInt const]
    {\seqq{type_{ce}(\sint{n}(c)) = (\mbox{``Gtyp''}, \sint{n})}}
   {\seqq{-}}
    \\[1ex]
    \inferii[mux]
    {type_{ce}(mux(e_0, e_1, e_2)) = maxTypes(t_1, t_2)}
   {\seqq{type_{ce}(e_1) = t_1}} {\seqq{type_{ce}(e_2) = t_2}}
    \quad
%    \\[1ex]
    \infer[validif]
    {\seqq{type_{ce}(validif(e_1, e_2)) = type_{ce}(e_2)}}
   {\seqq{-}}
    \quad
    \\[1ex]
    \inferii[$bop$=add, sub, $\theta=UInt, SInt$]
    {type_{ce}(bop(e_1, e_2)) = (\mbox{``Gtype''}, \theta\mbox{<}\max(n_1, n_2)+1\mbox{>})}
   {\seqq{type_{ce}(e_1) = (\mbox{``Gtype''}, \theta\mbox{<}n_1\mbox{>})}} {\seqq{type_{ce}(e_2) =  (\mbox{``Gtype''}, \theta\mbox{<}n_2\mbox{>})}}
    \quad
    \\[1ex]
    \inferii[mul, $\theta=UInt, SInt$]
    {type_{ce}(mul(e_1, e_2)) = (\mbox{``Gtype''}, \theta\mbox{<}n_1+ n_2\mbox{>})}
   {\seqq{type_{ce}(e_1) = (\mbox{``Gtype''}, \theta\mbox{<}n_1\mbox{>})}} {\seqq{type_{ce}(e_2) =  (\mbox{``Gtype''}, \theta\mbox{<}n_2\mbox{>})}}
    \quad
    \\[1ex]
    \inferii[div, UInt]
    {type_{ce}(div(e_1, e_2)) = (\mbox{``Gtype''}, \uint{n_1})}
   {\seqq{type_{ce}(e_1) = (\mbox{``Gtype''}, \uint{n_1})}} {\seqq{type_{ce}(e_2) =  (\mbox{``Gtype''}, \uint{n_2})}}
    \quad
    \\[1ex]
    \inferii[div, SInt]
    {type_{ce}(div(e_1, e_2)) = (\mbox{``Gtype''}, \sint{n_1+1})}
   {\seqq{type_{ce}(e_1) = (\mbox{``Gtype''}, \sint{n_1})}} {\seqq{type_{ce}(e_2) =  (\mbox{``Gtype''}, \sint{n_2})}}
    \quad
    \\[1ex]
  \end{gather*}
%  \vspace{-4mm}
  \caption{Definition of $type_{ce}(e)$}
  \label{tab:type-hexp}
\end{table}

$type_{ce}(e)$ is defined as follows.
\begin{itemize}
\item $type_{ce}(\uint{n}(c)) = (\mbox{``Gtyp''}, \mbox{``UInt''})$,
%
\item $type_{ce}(\sint{n}(c)) = (\mbox{``Gtyp''}, \mbox{``SInt''})$,
%
\item let $type_{ce}(e_1) = t_1$ and $type_{ce}(e_2) = t_2$, then $type_{ce}(mux(e_0, e_1, e_2)) = maxTypes(t_2, t_3)$,
%
\item $type_{ce}(validif(e_1, e_2)) = type_{ce}(e_2)$,
%
\item let $type_{ce}(e_1) =t_1$ and $type_{ce}(e_2) = t_2$, then
\begin{itemize}
\item $type_{ce}(add(e_1, e_2)) =\theta \bwidth{\max(n_1, n_2)+1}$ and $type_{ce}(sub(e_1, e_2)) = \theta \bwidth{\max(n_1, n_2)+1}$,
%
\item $type_{ce}(mul(e_1, e_2)) = \theta \bwidth{n_1+n_2}$,
%
\item $type_{ce}(div(e_1, e_2)) = UInt \bwidth{n_1}$ if $\theta = UInt$, and $type_{ce}(div(e_1, e_2)) = SInt \bwidth{n_1+1}$ otherwise,
%
\item $type_{ce}(mod(e_1, e_2)) = \theta \bwidth{\min(n_1,n_2)}$,
%
\item $type_{ce}(cp(e_1, e_2)) = UInt \bwidth{1}$, where $cp \in \{lt, leq, gt, geq, eq, neq\}$,
%
\item $type_{ce}(dshl(e_1, e_2)) = \theta \bwidth{n_1+2^{n_2}-1}$,
%
\item $type_{ce}(dshr(e_1, e_2)) = \theta \bwidth{n_1}$,
%
\item $type_{ce}(and(e_1, e_2)) = type_{ce}(or(e_1, e_2)) =  te(xor(e_1, e_2)) = \theta \bwidth{\max(n_1, n_2)}$,
%
\item $type_{ce}(cat(e_1, e_2)) = \theta \bwidth{n_1 + n_2}$,
\end{itemize}
%
\item let $te(e_1) = \theta \bwidth{n_1}$, where $\theta \in \{UInt, SInt\}$, then
\begin{itemize}
\item $te(andr(e_1)) = te(orr(e_1)) =  te(xorr(e_1)) = \uint{1}$,
%
\item $te(neg(e_1)) = \sint{n_1 +1}$, $te(not(e_1)) =  \uint{n_1}$,
%
\item $te(cvt(e_1)) = \sint{n_1 +1}$ if $\theta = UInt$, and $te(cvt(e_1)) =  \sint{n_1}$ otherwise,
\end{itemize}
%
\item let $te(e_1) = \theta \bwidth{n_1}$, where $\theta \in \{UInt, SInt\}$, then
\begin{itemize}
\item $te(pad(e_1, n)) = \uint{n}$,
%
\item $te(shl(e_1,n)) = \theta \bwidth{n_1 + n}$,
%
\item $te(shr(e_1, n)) = \theta \bwidth{n_1-n}$ if $n_1 > n$, and $te(shr(e_1, n)) =  \theta \bwidth{1}$ otherwise,
%
\item $te(head(e_1, n)) = \uint{n}$, $te(tail(e_1, n)) = \uint{n_1-n}$,
%
\item $te(bits(e_1, n, n')) = \uint{n'-n+1}$.
\end{itemize}
%
\end{itemize}
}
